\chapter{Relación con los Objetivos de Desarrollo Sostenible}\label{ap:ods}

\section{ODS 9: Industria, innovación e infraestructura}\label{sec:ods9}

El ODS~9 persigue construir infraestructuras resilientes, promover la industrialización
inclusiva y sostenible y fomentar la innovación. Este es el objetivo con el que el
trabajo guarda una relación más directa.

La memoria desarrolla y evalúa una arquitectura de aprendizaje profundo
(\emph{Stockformer}) y estudia su transferibilidad a un mercado distinto del original,
el S\&P~500. La aportación principal es de carácter metodológico.
Se proponen y comparan variantes del modelo, funciones de pérdida orientadas al
\emph{ranking} transversal y esquemas de validación \emph{walk-forward} que permiten
medir la estabilidad temporal del modelo.

Además, el trabajo se acompaña de una infraestructura de investigación reproducible:
un repositorio organizado, configuraciones versionadas de cada experimento, scripts de
construcción de datos y de evaluación, y una batería de pruebas automáticas. Esta
infraestructura es también una contribución alineada con el ODS~9, porque reduce
la barrera para que otros investigadores repliquen, auditen y extiendan los resultados.

\section{ODS 8: Trabajo decente y crecimiento económico}\label{sec:ods8}

El ODS~8 promueve el crecimiento económico sostenido e inclusivo y el buen
funcionamiento de los sistemas económicos. La relación del trabajo con este objetivo se
articula a través de la \emph{transparencia} en la evaluación del rendimiento.

Una parte sustancial de la memoria se dedica a la atribución del retorno \emph{neto}: a
descomponer el rendimiento de la estrategia descontando costes de transacción,
rotación y exposición a factores de riesgo. Este énfasis en medir el resultado real,
y no únicamente el bruto, contribuye a una cultura de evaluación más honesta y prudente
en el ámbito de la inversión cuantitativa, donde es frecuente sobreestimar el
rendimiento al ignorar fricciones de mercado.

La contribución al ODS~8 no consiste en garantizar rentabilidades ni en proponer un
producto de inversión. De hecho, el trabajo evita ese tipo de afirmaciones: sus
conclusiones son metodológicas y sus resultados se presentan con intervalos de
confianza y contrastes de significación estadística que delimitan lo que puede y no
puede afirmarse. Esta prudencia es importante en un contexto financiero.

\section{ODS 12: Producción y consumo responsables}\label{sec:ods12}

El ODS~12 aboga por modelos de producción y consumo responsables, lo que en el contexto
del aprendizaje profundo se traduce en un uso eficiente de los recursos
computacionales.

El modelo estudiado emplea mecanismos de atención espacial dispersa y convoluciones
temporales que reducen el coste de cómputo respecto a una atención densa sobre todos
los pares de activos. La memoria, además, incorpora explícitamente el tiempo de cómputo
como una de las métricas de comparación entre modelos (véase la escalera de complejidad
del Apéndice~\ref{ap:tablas}), lo que permite valorar cada mejora de precisión frente a
su coste energético y computacional.

Esta perspectiva fomenta decisiones de diseño responsables. En igualdad de
condiciones, se prefieren los modelos más simples y eficientes, y el gasto computacional
adicional se justifica únicamente cuando aporta una mejora medible y estadísticamente
significativa.
La contribución al ODS~12 es, por tanto, la de promover una práctica investigadora
consciente del coste de los recursos, más que la de un impacto material directo.